\documentclass{article}
\usepackage[margin=2cm]{geometry}
\usepackage{amsmath,amssymb}
\usepackage{booktabs}
\usepackage{array}
\usepackage{enumitem}
\usepackage[most]{tcolorbox}

% Comando para isolar as referências alinhadas à direita no título de cada passo
\newcommand{\stepref}[1]{\hfill\mbox{\footnotesize\normalfont\textsf{\cite{biswas2017}, #1}}}

\newtcolorbox{algorithmbox}[1][]{%
    enhanced,
    breakable,
    colback=white,
    colframe=black,
    fonttitle=\bfseries,
    coltitle=black,
    attach boxed title to top left={xshift=4mm, yshift=-2mm},
    boxed title style={colback=white, colframe=black, boxrule=0.5pt},
    boxrule=0.8pt,
    arc=0mm,
    top=5mm,
    bottom=3mm,
    left=3mm,
    right=3mm,
    title=#1
}

\begin{document}

\section*{Mapeamento de Vari\'aveis}

A Tabela~\ref{tab:mapa_variaveis} correlaciona a nota\c{c}\~ao adotada no algoritmo com os s\'imbolos originais do artigo de Biswas et al.~\cite{biswas2017}.

\begin{table}[htbp]
\centering
\small
\caption{Mapeamento, dimens\~ao e correspond\^encia das vari\'aveis do SPUKF.}
\label{tab:mapa_variaveis}
\begin{tabular}{llcp{6.8cm}l}
\toprule
\textbf{S\'imbolo} & \textbf{Biswas} & \textbf{Dimens\~ao} & \textbf{Descri\c{c}\~ao no Filtro de Navega\c{c}\~ao} & \textbf{Ref. Biswas} \\
\midrule
$p$ & $n$ & Escalar & Dimens\~ao do vetor de estado ($p = 15$) & Sec.~II \\
$\hat{x}_{k-1|k-1}$ & $\hat{Y}^+(k)$ & $p \times 1$ & Estimativa m\'edia a posteriori no instante $t_{k-1}$ & eq.~(4) \\
$P_{k-1|k-1}$ & $P_{YY}^+(k)$ & $p \times p$ & Covari\^ancia do erro a posteriori no instante $t_{k-1}$ & eq.~(5) \\
$\Delta x_i$ & $\Delta Y_i$ & $p \times 1$ & Perturba\c{c}\~ao determin\'istica do $i$-\'esimo ponto sigma & eq.~(5) \\
$\kappa$ & $\kappa$ & Escalar & Par\^ametro de escala da transformada unscented & eq.~(5) \\
$W_i$ & $W_i$ & Escalar & Pesos escalares da transformada para m\'edia e covari\^ancia & eq.~(5) \\
$u_{k-1}$ & --- & $6 \times 1$ & Entrada inercial: for\c{c}a espec\'ifica e velocidade angular & --- \\
$\Delta t$ & $\delta t$ & Escalar & Intervalo temporal de propaga\c{c}\~ao ($t_k - t_{k-1}$) & eq.~(15) \\
$f(x, u)$ & $f(t, Y, \nu)$ & $p \times 1$ & Modelo cont\'inuo da din\^amica cinem\'atica & eq.~(1) \\
$x_0(t_k)$ & $Y_0(t+\delta t)$ & $p \times 1$ & Trajet\'oria central integrada numericamente & eqs.~(15)--(16) \\
$\mathcal{J}$ & $\mathcal{J}$ & $p \times p$ & Jacobiana cont\'inua avaliada em $\hat{x}_{k-1|k-1}$ & eq.~(22) \\
$\Phi$ & $\frac{\partial F}{\partial Y}$ & $p \times p$ & Matriz de transi\c{c}\~ao discreta ($\Phi = \mathrm{e}^{\mathcal{J}\Delta t}$) & eq.~(22) \\
$\tilde{x}_i(t_k)$ & $\tilde{Y}_i(t+\delta t)$ & $p \times 1$ & Pontos sigma mapeados por expans\~ao de 1.\textordfeminine\ ordem & eq.~(21) \\
$\hat{x}_{k|k-1}$ & $\tilde{Y}^-(t+\delta t)$ & $p \times 1$ & Estimativa de estado m\'edia a priori no instante $t_k$ & eq.~(25) \\
$P_{k|k-1}$ & $\tilde{P}_{YY}^-(t+\delta t)$ & $p \times p$ & Covari\^ancia do erro do estado a priori & eq.~(28) \\
$h(\cdot)$ & $h(\cdot)$ & $m \times 1$ & Modelo de observa\c{c}\~ao GNSS ($m = 3$) & eq.~(2) \\
$\mathcal{Y}_i$ & $\tilde{Z}_i$ & $m \times 1$ & Ponto sigma transformado no espa\c{c}o de medi\c{c}\~ao & eq.~(31) \\
$\bar{y}_k$ & $\tilde{Z}^-$ & $m \times 1$ & Vetor de medi\c{c}\~ao predito a priori & eq.~(32) \\
$R$ & $R$ & $m \times m$ & Covari\^ancia do ru\'ido de medi\c{c}\~ao GNSS & eq.~(10) \\
$P_{hh}$ & $\tilde{S}$ & $m \times m$ & Matriz de covari\^ancia da inova\c{c}\~ao residual & eq.~(35) \\
$P_{gh}$ & $\tilde{P}_{YZ}$ & $p \times m$ & Matriz de covari\^ancia cruzada estado--medi\c{c}\~ao & eq.~(34) \\
$K$ & $\tilde{K}$ & $p \times m$ & Matriz de ganho de Kalman & eqs.~(12), (36) \\
$y_k$ & $Z(k+1)$ & $m \times 1$ & Vetor de medi\c{c}\~ao recebido pelo recetor GNSS & eq.~(12) \\
$\hat{x}_{k|k}$ & $\hat{Y}^+(k+1)$ & $p \times 1$ & Estimativa de estado corrigida a posteriori & eqs.~(12), (37) \\
$P_{k|k}$ & $P_{YY}^+(k+1)$ & $p \times p$ & Covari\^ancia do erro atualizada a posteriori & eqs.~(12), (38) \\
\bottomrule
\end{tabular}
\end{table}

\vspace{0.3cm}
\newpage

\begin{algorithmbox}[Quadro 3.1: Algoritmo SPUKF (Biswas et al.) com notação alinhada]
\begin{enumerate}[leftmargin=*, itemsep=6pt]
    \item \textbf{Gera\c{c}\~ao dos Pontos Sigma no Instante $t_{k-1}$:} \stepref{eqs.~(4)--(5)} \\
    A partir da estimativa a posteriori $\hat{x}_{k-1|k-1} \in \mathbb{R}^p$ e da covari\^ancia $P_{k-1|k-1} \in \mathbb{R}^{p \times p}$, gera-se o conjunto de $2p+1$ pontos sigma:
    \begin{align}
        x_0(t_{k-1}) &= \hat{x}_{k-1|k-1}, \\
        x_i(t_{k-1}) &= \hat{x}_{k-1|k-1} + \Delta x_i, \quad i = 1, \dots, 2p,
    \end{align}
    onde as perturba\c{c}\~oes $\Delta x_i$ s\~ao extra\'idas das colunas da decomposi\c{c}\~ao matricial:
    \begin{equation}
        \Delta x_i = 
        \begin{cases} 
            \left(\sqrt{(p+\kappa)P_{k-1|k-1}}\right)_i, & i = 1, \dots, p, \\ 
            -\left(\sqrt{(p+\kappa)P_{k-1|k-1}}\right)_{i-p}, & i = p+1, \dots, 2p,
        \end{cases}
    \end{equation}
    com os pesos escalares associados:
    \begin{equation}
        W_0 = \frac{\kappa}{p+\kappa}, \quad 
        W_i = \frac{1}{2(p+\kappa)}, \quad i = 1, \dots, 2p.
    \end{equation}

    \item \textbf{Propaga\c{c}\~ao \'Unica e Mapeamento de 1.\textordfeminine\ Ordem:} \stepref{eqs.~(15)--(16), (20)--(22)} \\
    Propaga-se numericamente apenas o estado central $x_0(t_{k-1}) = \hat{x}_{k-1|k-1}$ atrav\'es da din\^amica cont\'inua $\dot{x}(t) = f(x(t), u(t))$ ao longo do intervalo $\Delta t$:
    \begin{equation}
        x_0(t_k) = F\left(t_{k-1}, \hat{x}_{k-1|k-1}, u_{k-1}\right) = \hat{x}_{k-1|k-1} + \int_{t_{k-1}}^{t_k} f(\tau, x(\tau), u(\tau)) \, \mathrm{d}\tau.
    \end{equation}
    Avalia-se a Jacobiana cont\'inua $\mathcal{J} = \left. \frac{\partial f}{\partial x} \right|_{\hat{x}_{k-1|k-1}}$, calcula-se a matriz de transi\c{c}\~ao $\Phi = \mathrm{e}^{\mathcal{J}\Delta t}$ e mapeiam-se os restantes $2p$ pontos sigma:
    \begin{equation}
        \tilde{x}_i(t_k) = x_0(t_k) + \mathrm{e}^{\mathcal{J}\Delta t} \Delta x_i, \quad i = 1, \dots, 2p.
    \end{equation}

    \item \textbf{C\'alculo da M\'edia Predita do Estado (A Priori):} \stepref{eq.~(25)} \\
    Determina-se a estimativa m\'edia a priori:
    \begin{equation}
        \hat{x}_{k|k-1} = \sum_{i=0}^{2p} W_i \tilde{x}_i(t_k).
    \end{equation}

    \item \textbf{C\'alculo da Covari\^ancia do Erro do Estado (A Priori):} \stepref{eq.~(28)} \\
    Calcula-se a matriz de covari\^ancia do erro a priori:
    \begin{equation}
        P_{k|k-1} = \sum_{i=0}^{2p} W_i \left[\tilde{x}_i(t_k) - \hat{x}_{k|k-1}\right] \left[\tilde{x}_i(t_k) - \hat{x}_{k|k-1}\right]^T.
    \end{equation}

    \item \textbf{Instancia\c{c}\~ao das Observa\c{c}\~oes:} \stepref{eq.~(31)} \\
    Instancia-se cada ponto sigma no modelo de medi\c{c}\~ao $h(\cdot)$:
    \begin{equation}
        \mathcal{Y}_i = h\left(\tilde{x}_i(t_k)\right), \quad i = 0, \dots, 2p.
    \end{equation}

    \item \textbf{C\'alculo da Medi\c{c}\~ao Predita:} \stepref{eq.~(32)} \\
    Calcula-se o vetor de observa\c{c}\~ao m\'edio predito:
    \begin{equation}
        \bar{y}_k = \sum_{i=0}^{2p} W_i \mathcal{Y}_i.
    \end{equation}

    \item \textbf{C\'alculo da Covari\^ancia da Inova\c{c}\~ao:} \stepref{eq.~(35)} \\
    Com a matriz de covari\^ancia do ru\'ido de medi\c{c}\~ao $R$, computa-se:
    \begin{equation}
        P_{hh} = \sum_{i=0}^{2p} W_i \left[\mathcal{Y}_i - \bar{y}_k\right] \left[\mathcal{Y}_i - \bar{y}_k\right]^T + R.
    \end{equation}

    \item \textbf{C\'alculo da Covari\^ancia Cruzada Estado--Medi\c{c}\~ao:} \stepref{eq.~(34)} \\
    Determina-se a matriz de correla\c{c}\~ao cruzada:
    \begin{equation}
        P_{gh} = \sum_{i=0}^{2p} W_i \left[\tilde{x}_i(t_k) - \hat{x}_{k|k-1}\right] \left[\mathcal{Y}_i - \bar{y}_k\right]^T.
    \end{equation}

    \item \textbf{Atualiza\c{c}\~ao da Medi\c{c}\~ao (A Posteriori):} \stepref{eqs.~(12), (36)--(38)} \\
    Determina-se o ganho de Kalman $K$ e atualizam-se a estimativa de estado e a covari\^ancia do erro perante a medi\c{c}\~ao $y_k$:
    \begin{align}
        K &= P_{gh} P_{hh}^{-1}, \\
        \hat{x}_{k|k} &= \hat{x}_{k|k-1} + K \left(y_k - \bar{y}_k\right), \\
        P_{k|k} &= P_{k|k-1} - K P_{hh} K^T.
    \end{align}
\end{enumerate}
\end{algorithmbox}

\begin{thebibliography}{1}

\bibitem{biswas2017}
S.~K. Biswas, L.~Qiao, and A.~G. Dempster, ``A novel a priori state computation strategy for the unscented Kalman filter to improve computational efficiency,'' \emph{IEEE Transactions on Automatic Control}, vol.~62, no.~4, pp.~1852--1864, Apr. 2017.

\end{thebibliography}

\end{document}